\documentclass[11pt,letterpaper]{article}
\usepackage[T1]{fontenc}
\usepackage{lmodern}
\usepackage[margin=0.88in]{geometry}
\usepackage{amsmath,amssymb,amsthm,mathtools}
\usepackage{booktabs,longtable,array,tabularx,microtype}
\usepackage{xcolor,xurl,enumitem,fancyhdr,needspace,etoolbox}
\usepackage[hidelinks,pdftitle={PA195: a 195-state binary Turing machine for the consistency of Peano arithmetic},
 pdfauthor={Research prepared for Vladimir Reshetnikov}]{hyperref}
\definecolor{ink}{HTML}{17364D}
\definecolor{accent}{HTML}{0B6B6A}
\newtheorem{theorem}{Theorem}[section]
\newtheorem{lemma}[theorem]{Lemma}
\newtheorem{proposition}[theorem]{Proposition}
\newtheorem{corollary}[theorem]{Corollary}
\newcommand{\PA}{\mathrm{PA}}
\newcommand{\Con}{\operatorname{Con}}
\newcommand{\Halt}{\operatorname{Halt}}
\newcommand{\file}[1]{\texttt{\nolinkurl{#1}}}
\newcommand{\code}[1]{\texttt{#1}}
\pretocmd{\section}{\Needspace{8\baselineskip}}{}{}
\AtBeginEnvironment{lemma}{\Needspace{5\baselineskip}}
\setlength{\emergencystretch}{2em}
\setlength{\parskip}{0.35em}
\setlength{\parindent}{0pt}
\setlist{itemsep=0.2em,topsep=0.4em}
\pagestyle{fancy}
\fancyhf{}
\fancyhead[L]{\small\color{ink}PA195}
\fancyhead[R]{\small A binary machine for $\Con(\PA)$}
\fancyfoot[C]{\thepage}
\setlength{\headheight}{14pt}
\title{\color{ink}\bfseries PA195\\[0.3em]
\Large A 195-state binary Turing machine\\for the consistency of Peano arithmetic}
\author{Research prepared for Vladimir Reshetnikov}
\date{13 November 2026\\[0.3em]\small Source snapshot:
\texttt{906367250a7f96c923147b10df392294bced7be5}}
\begin{document}
\maketitle
\thispagestyle{empty}
\begin{abstract}
We give an explicit deterministic single-tape, two-symbol Turing machine
with 195 working states, started on a two-way infinite all-zero tape, and
an ordinary mathematical proof that Peano arithmetic proves its
nonhalting equivalent to $\Con(\PA)$. The construction preserves the
191-item diagonal-vacuity axiom emitter in the supplied research corpus.
A new binary association admits seven simultaneously valid
history-sensitive reader identifications: 134 inherited controls plus
68 reader controls, less seven, gives 195.

The packet contains the complete 390-row table, deterministic
reconstruction, all 3,700 source instructions, 28 finite local simulation
certificates, an exact 58,154-step startup certificate, and executable
symbolic checks for all 4,815 primitive branches. An independent audit
checks the source graph, two representations of startup, all protected
engine rows, and 19,260 concrete instruction frames.
The comparison covers all 27 incoming archives, the direct PA research
directory, 20 designated U208 graph classes, and 25 distinct supported
196-state layouts. The new arithmetic endpoint is proved in ordinary
mathematical form; a new end-to-end Lean or Rocq theorem is not claimed.
\end{abstract}
\vspace{0.6em}
\begin{center}
\fcolorbox{accent}{white}{\parbox{0.88\linewidth}{\centering
\large\textbf{195 working states \quad 2 symbols \quad 390 transitions}\\[0.4em]
\(\displaystyle
 \PA\vdash
 \neg\Halt(\mathrm{PA195},\mathrm{blank})\ \leftrightarrow\ \Con(\PA).
\)}}
\end{center}
\newpage
\tableofcontents
\newpage

\section{Result and exact machine conventions}
\label{sec:result}
The deliverable \file{pa195/pa195.tm} is the definitive literal machine.
Its working labels are $0,\ldots,194$, its start label is $0$, and its
separate halt label is $195$. The tape alphabet is $\{0,1\}$, the tape
extends without bound in both directions, the initial head coordinate is
zero, and every initial cell is zero. There is one transition for every
working-state/read-symbol pair. No stay-put move is used.

A row \((q,b,w,D,q')\) means: read \(b\), write \(w\), move one cell in
direction \(D\in\{L,R\}\), and enter \(q'\). The halt control has no
outgoing row and is excluded from the working-state count. All 195
working states are reachable in the finite transition graph. This last
claim concerns graph reachability, not occurrence on the actual blank run.
\begin{center}
\begin{tabular}{@{}ll@{}}
\toprule
Object & Exact value\\
\midrule
Working states & 195\\
Transition rows & 390\\
Start / halt & 0 / 195\\
Source root in the dense table & 75\\
Unique halt transition & \(63,\ 1\ \mapsto\ 1,\ R,\ 195\)\\
First transition on blank & \(0,\ 0\ \mapsto\ 1,\ R,\ 1\)\\
Source instructions & 3,700\\
Primitive branch contracts & 4,815\\
Exact startup length & 58,154 literal transitions\\
\bottomrule
\end{tabular}
\end{center}
The SHA-256 of the literal five-column table, including its final newline, is
\begin{quote}\small
\nolinkurl{5c6afd87f7c4f24d9ad8e9e302cc21df14b6f2e3715d03322a75463a431d8455}.
\end{quote}
Appendix~\ref{app:table} prints the entire table in paired-column form.
The native-label version and its complete conjugacy are also included;
native label 254 and dense label 195 denote the same halt control.

\subsection{What is new}
The new result is a physical compilation of the existing
\emph{diagonal191} program below the 196-state threshold. Its theory,
formula encoding, axiom order, formula parameters, register assignment,
and exact ordered instruction sequence are unchanged from the
\file{diagonal196} construction in the supplied archive \cite{diagonal}.
The source parenthesization and reader allocation change.
The ordinary PA proof in Section~\ref{sec:pa-calibration} explicitly
connects this unchanged calculus to the delivered table.

For the supplied \file{diagonal196} parent, the allocation is
\(134+68-6=196\).
For the new association, it is
\[
 \boxed{134+68-7=195.}
\]
The new association therefore retains the same number of allocated
readers before folding but admits one additional \emph{compatible}
fold. Optimizing only the number of syntax nodes would miss this saving.
The 134 inherited controls comprise 36 engine controls, 13 leaf routers,
and 85 readers outside this AX allocation. The final reader census is
$85+(68-7)=146$; hence $36+13+146=195$.

\subsection{Meaning of the proof claim}
The theorem is the internal PA biconditional, not merely external
nonhalting under the standard interpretation. It is proved here by
uniform syntactic proof transformations, arithmetic induction, and
positive finite simulation. Its validity does not assume $\Con(\PA)$.
Actual blank-tape nonhalting follows if PA is consistent.

The published PA196 Lean development supplies the internal implication
from nonhalting to consistency, together with its other external and
independence results \cite{chang,bridge}. It should not be described as
already containing the converse internal implication. This report
supplies the converse by syntactic PA proof annotations and finite
bit-comprehension proof generation. The executable machine certificates
and the ordinary proof are available for review; a newly elaborated
Lean or Rocq endpoint for PA195 remains future formalization work.

\section{Comparison of the supplied designs}
\label{sec:comparison}
\subsection{Pinned scope and deduplication}
The comparison fixes ProveIt commit
\file{906367250a7f96c923147b10df392294bced7be5}.
The direct \file{Computability/BusyBeaver/Research/PA} directory contains
the imported PA196 and PA209 projects plus their comparison and source
experiments \cite{comparison}. The incoming directory has 27 ZIP
archives, totaling 224,873,862 bytes, with 21,003 non-directory outer
members. Text, code, and machine members were inspected or indexed;
nested archives were inspected to depth two. Appendix~\ref{app:archives}
gives a contribution/status entry for every incoming archive.

The numeric table census finds 1,857 \file{.tm} occurrences and 521
distinct literal byte strings. These are \emph{file} counts, not claims
of 521 correct consistency machines: they include repeated ancestors,
experimental allocations, rejected candidates, and small test controls.
For meaningful comparisons we canonicalize designated tables as rooted,
read-symbol- and action-preserving graphs. A canonicalization preserves
the start control, the halt control, writes, directions, and the two
read-symbol labels; it does not collapse a change of tape convention.

The designated U208 corpus has 27 presentations and 20 graph classes.
The first nine presentations form two classes, with five and four
presentations respectively. Eighteen later helper variants each add a
class. Twenty-five distinct supported 196-state layouts are retained.
Twenty-three have complete evidence in the supplied reports; two more
have their missing startup or contract checks completed in this audit.
\emph{Supported} means that the table has its stated source, runtime,
startup, and mathematical evidence.
It does not upgrade every report to a kernel-checked theorem.

The full 196-state census contains 268 table occurrences, 131 distinct
byte strings, and 90 rooted graph classes. Every class is classified in
the accompanying ledger; the remaining classes have partial, bounded,
or negative evidence and are not silently promoted to the supported frontier.

The incoming corpus also contains 14 occurrences of 195-state tables:
eight literal byte strings and six candidate constructions after
native/dense renamings are identified. All are rejected or uncertified
in their source reports. Some fail an exact primitive contract; others
have a wrong source successor or an explicit blank-tape halt. They are
not predecessors already meeting the requested endpoint.

\clearpage
\subsection{The 196-state frontier}
{\small
\setlength{\tabcolsep}{5pt}
\begin{longtable}{@{}p{0.34\linewidth}rrrrl@{}}
\toprule
Design & States & Sites & Contracts & Startup & Evidence\\
\midrule
\endfirsthead
\toprule
Design & States & Sites & Contracts & Startup & Evidence\\
\midrule
\endhead
\midrule
\endfoot
\bottomrule
\endlastfoot
Published PA196 & 196 & 3,812 & 4,959 & 58,154 & K\\
Alternative head order & 196 & 3,812 & 4,959 & 58,154 & E\\
Closing CD3 & 196 & 3,812 & 4,959 & 58,154 & I\\
Closing CD4 & 196 & 3,812 & 4,959 & 58,154 & I\\
Closing CD5 & 196 & 3,812 & 4,959 & 58,154 & I\\
Closing CD6 & 196 & 3,812 & 4,959 & 58,154 & I\\
Rotation layout 1 & 196 & 3,812 & 4,959 & 58,294 & E\\
Rotation layout 2 & 196 & 3,812 & 4,959 & 58,294 & E\\
Eight-zero scanner layout & 196 & 3,812 & 4,959 & 59,164 & E\\
A/D register swap & 196 & 3,812 & 4,959 & 58,122 & E\\
Rebracketed source & 196 & 3,812 & 4,959 & 58,154 & E\\
H-only Shift & 196 & 3,811 & 4,958 & 57,922 & E\\
Twelve-rotation layout & 196 & 3,812 & 4,959 & 58,294 & E\\
Combined196 & 196 & 3,788 & 4,928 & 57,922 & E\\
Diagonal196 & 196 & 3,700 & 4,815 & 58,154 & E+L\\
Separate196 & 196 & 3,673 & 4,780 & 58,154 & E+L\\
Separate/scanner196 & 196 & 3,673 & 4,780 & 58,294 & E+L\\
Original/scanner196 & 196 & 3,812 & 4,959 & 58,294 & E\\
Zerojump196 & 196 & 3,812 & 4,959 & 58,402 & I\\
Diagonal closing CD3 & 196 & 3,700 & 4,815$^{*}$ & 58,154 & I+L\\
Diagonal closing CD4 & 196 & 3,700 & 4,815$^{*}$ & 58,154 & I+L\\
Diagonal closing CD5 & 196 & 3,700 & 4,815$^{*}$ & 58,154 & I+L\\
Diagonal closing CD6 & 196 & 3,700 & 4,815$^{*}$ & 58,154 & I+L\\
Zerojump CD3 & 196 & 3,812 & 4,959 & 58,366 & I$^{\dagger}$\\
Rebracketed Combined196 & 196 & 3,788 & 4,928 & 57,922 & E$^{\dagger}$\\
\midrule
\textbf{New PA195} & 195 & 3,700 & 4,815 & 58,154 & E+L\\
\end{longtable}
}

\textbf{Evidence key.} K: the published Lean development, including its
internal implication from nonhalting to consistency. E: complete executable
symbolic runtime and exact startup, with ordinary mathematical transfer.
I: runtime uses a proved source invariant. L: the numerical calculus changes
and therefore needs its stated logical calibration. None of E, I, or L
denotes a newly compiled kernel theorem. The new PA195 row has a full
ordinary proof of the internal biconditional in this report.

$^{*}$For each diagonal closing variant, 45 affected branches are replayed
as 69 invariant-restricted affine cases; every other baseline branch avoids
the changed row. This differential proof covers all 4,815 baseline branches.

$^{\dagger}$The CD3 zerojump and rebracketed Combined196 tables are archived
designs whose missing complete contract or startup check was supplied in
this audit. Their fresh receipts are included.

The source counts are primitive instruction occurrences. Contract counts
include both outcomes of a conditional decrement. Prior numerical
figures are supported by the supplied receipts, with the two audit-completed
cases marked separately; the PA195 row is newly recomputed. Exact archive/member/hash bindings are in
\file{comparison/frontier196_and_pa195.json}, and complete state
bijections for the 25 normalized tables are retained under
\file{comparison/prior196/}.

\subsection{Why the older reductions are not additive}
James Chang's specified 209-state predecessor and the U209 parents of the
U208 archives are distinct machines. The U208 saving identifies the
particular native roles 172 and 211 using entry history and an existing
helper. Those old roles are already absent from PA196. One cannot simply
subtract that saving again from a different 196-state table.

The published 209-to-196 change combines the transposed quantifier code,
a smaller equality interface, changed axiom order, and contextual reader
sharing \cite{chang}. The original PA196 has 3,812 source sites and 159
distinct internal source nodes; its 201-item AX word uses 68 allocated
reader nodes and six folds. The diagonal-vacuity source reduces the AX
word to 191 and the full source to 3,700 sites. Its earlier physical
layouts reached 196, and provide the parent used here.

Several apparently simpler proposals do not yield a smaller binary
machine. A 157-node source can lose a required helper. A different
158-node layout needs an extra scanner state. The five-zero source
representation changes the initial source address and loses inherited
contextual sharing. The separate-vacuity emitter has only 190 macro
items, but pushes ten stack entries instead of nine and has a different
command compiler interface. These examples explain why the present
search scores an actual compatible allocation.

The direct PA directory's additional literal tables are inventoried in
\file{comparison/pa_directory_inventory.json}. Besides the two
published tables and alternative196, these include 197-state
negation/MP and shortcut experiments, source-pruning allocations, an
identity reconstruction, and a 198-state root-frame experiment with an
unsettled blank initialization. None supplies a smaller accepted table.

\subsection{Ideas reused from ZFC and ConciseIPRA}
The useful ZFC inheritance is the compact unary-register engine and its
symbolic block contracts, source-tree sharing, and history-sensitive
reader reuse \cite{zfc}. The new machine keeps this engine and changes
its compiled proof-search program representation.

The useful proof-theoretic idea is to derive only the equality interface
actually used by the refutation compiler, exploiting a formula code as a
fresh variable label. The diagonal version is already in the incoming PA
research, and is credited as such. A small iPRA checker cannot simply
replace the searched calculus: the PA translation uses induction for an
arbitrary formula in finite bit comprehension. Section~\ref{sec:pa-calibration}
explains this strength distinction; no identification of
$\Con(\mathrm{iPRA})$ with $\Con(\PA)$ is used \cite{ipra}.

\section{Exact source and compilation}
\label{sec:source}
\subsection{Source representation and invariant interfaces}
The delivered \file{pa195/source.json} specifies every primitive source
instruction, its binary address, and every ordinal successor. The
instruction vocabulary is increment, conditional decrement, break,
continue, and halt. Binary association determines how break and
continue depths are resolved; equal source fragments can share physical
reader states. Consequently equality of flattened opcodes alone would
not be sufficient. The verifier compares the complete ordinal
control-flow graphs with two independently structured resolvers.

The 3,700 sites consist of 1,468 increments, 1,115 conditional
decrements, 1,116 jumps, and one halt. Thus the branch census is
\[
 1468+2(1115)+1116+1=4815.
\]
The longest path has 37 bits and the longest consecutive zero run is six.
The clean seed remains source ordinal 25, at path \code{100000}.

At a clean physical frame the left marker is immediately left of the
first path cell, a zero-filled gap separates the path from the right
header, and the complete register file follows the header. A natural
value \(v\) is encoded by \(0\,1^{v+1}\).
Logical registers \(0,\ldots,8\) occupy physical positions
\[
 (6,2,0,4,3,7,8,5,1).
\]
The symbolic gap is \(48+n\) for arbitrary \(n\ge0\).
Increment preserves the gap and appends a one-valued spare block;
a positive decrement increases the gap by one; zero decrement and
jump preserve it. The proof covers every complete finite file of at
least ten registers satisfying the indicated branch condition.

\subsection{The emitted numerical calculus}
Section~\ref{sec:pa-calibration} defines the formula constructors and
the changed gate fully. For completeness, the six unchanged heads are
listed here. \(D,A,H\) are arbitrary numerical parameters, read as
variable labels or formula codes in the indicated positions.
\begin{center}\small
\begin{tabular}{@{}cl@{}}
\toprule
Head & Formula\\
\midrule
1 & \((\exists x_D\,U(D,H))\to H\)\\
2 & \(U(D,U(A,H))\to U(A,U(D,H))\)\\
3 & \(M(H,D)\to(E(H,A)\to M(A,D))\)\\
4 & \(E(H,D)\to(M(A,H)\to M(A,D))\)\\
5 & \(\exists x_D\,E(H,D)\)\\
6 & \(U(D,A)\to(U(D,A\to H)\to U(D,H))\)\\
\bottomrule
\end{tabular}
\end{center}
The restricted propositional schema is
\[
 \rho_8(D,A,H,u)=((H\to u)\to D)\to((D\to H)\to(A\to H)),
 \qquad u\le D.
\]
The only primitive inference rules are modus ponens and generalization
of an already derived theorem. There is no extra equality rule or
arbitrary-tautology oracle.

In a command store \((t,w,d,h)\), an AX call with parameters \((c,x,y)\)
uses \(D=d+x\), \(H=h+y\), \(A=c\). In execution order it pushes
the saved \(t\), \(U(d,h)\), \(\rho_8(D,A,H,d)\), and heads 1 through 6.
Its new top is the diagonal gate when the guard succeeds and truth
otherwise; \(d,h\) are cleared. The exact output, including its full
stack order and the gate, is verified as an identity in a free algebra.
It is not inferred from selected small numerical codes.

The new \file{grammar.json} defines the complete association of the same
191 macro tokens. A pair of adjacent increments uses an already
available macro root, so the grammar's number of named rules must not
be confused with its 68 newly allocated reader controls. The deterministic
builder expands the exact grammar, retains the inherited macro entry
controls, and then applies the seven folds.

\subsection{Exact startup and total primitive simulation}
After exactly 58,154 literal steps, the delivered machine is in dense
state 75 (native 84), with head at absolute coordinate \(-2\).
Relative to that head the left marker is at \(-1\), the path is
\code{100000}, the right header is at 48, and the physical register
file is
\[
 [0,0,0,0,0,0,2,0,0]\ \mathbin{+\!\!+}\ [1]^{26}.
\]
The full finite nonzero support, containing 66 cells, is recorded in
\file{startup_certificate.json}; every other tape cell is zero.
The proof therefore uses blank initialization of this exact machine,
not an externally supplied nonblank input.

For both the unfolded and folded tables the executable symbolic
interpreter checks all 4,815 branch contracts at arbitrary gap and
register values. Its register oracle is justified by the inherited
unary-loop proofs. Independently, all 44 rows named in the Lean
\code{boxRows} contract, the additional increment-tail row, and even
the full 46-row protected superset were compared literally.
Every row is unchanged. The runtime soundness argument and its
PA internalization are developed in Section~\ref{sec:pa-calibration}.

% Standalone section fragment for the accompanying report.
\section{The seven-state quotient: a global finite certificate}

Let $R$ be the explicitly supplied unfolded 202-state table and let $M$
be \texttt{pa195.native.tm}.  Both use native start label $74$ and native
halt label $254$.  The dense presentation of $M$ has start $0$, working
labels $0,\ldots,194$, and halt label $195$.
The quotient map $\sigma$ identifies the following seven pairs and fixes
all other controls.  The first member of each pair is removed.
\[
\begin{gathered}
(2018,2024),\ (2030,2050),\ (2038,2019),\ (2042,2028),\\
(2053,2043),\ (2055,2034),\ (2062,2036).
\end{gathered}
\]
The unfolded machine consists of 134 inherited controls and 68 newly
allocated AX reader controls.  All seven identifications are disjoint, so
$202-7=195$ working controls remain.  The root-reachability census visits
every one of the 195 controls; no distinguished halt control is counted.

\paragraph{The incoming-history invariant.}
For every pair $(a,b)$ in this list, every incoming row of $R$ whose target
is $a$ writes $0$ and moves right; every incoming row whose target is $b$
writes $1$ and moves right.  There are exactly 22 such incoming rows, all
listed in \texttt{local\_fold\_certificate.json}.  Since neither $a$ nor
$b$ is the start control, every reachable configuration of $R$ satisfies
\[
 q=a\ \Longrightarrow\ t(h-1)=0,
 \qquad
 q=b\ \Longrightarrow\ t(h-1)=1,
\]
where $h$ is the head position and $t$ is the tape.  This is a direct
one-step induction.  It applies to every initial tape, not only to the
blank run or to well-formed register-program frames.

\paragraph{The 28 local obligations.}
For each of the 14 old controls and each current bit $v\in\{0,1\}$,
start $M$ in control $\sigma(q)$ at head position $0$, with the prescribed
history bit at $-1$ and current bit $v$ at $0$.
The certificate gives a positive finite run that reproduces the single
$R$ transition, including its write, head displacement, and projected
successor.  Until that endpoint every read and write is at $-1$ or $0$.
The two-cell tape at the endpoint is compared exactly.  Hence every
extension of those two cells to a bi-infinite tape satisfies the same
identity.  Of the 28 runs, ten have length 1, seventeen have length 3,
and one has length 5.  All other old transitions commute with $\sigma$
in exactly one step.

\paragraph{Halting preservation in both directions.}
Map a configuration $(q,h,t)$ to $(\sigma(q),h,t)$.  The preceding finite
identities and incoming-history invariant show that every step of $R$
from a reachable configuration is simulated by one to five steps of $M$.
Every intermediate control has an outgoing transition unless the final
endpoint is the halt control; $\sigma$ fixes that halt control.
Concatenation therefore sends any $R$ halting trace to an $M$ halting
trace.

Conversely, suppose $M$ halts after $T$ transitions.  If $R$ does not halt
within its first $T$ transitions, its first $T+1$ steps exist.  Their
simulating blocks concatenate to more than $T$ defined $M$ transitions,
in contradiction with determinism and the absence of outgoing rows at
halt.  Consequently $R$ halts within $T$ steps.  This also covers halts
inside an attempted comparison block; an intermediate premature halt
would contradict its verified finite trace.

All quantified tape and trace operations in this proof have standard
primitive-recursive finite encodings.  The invariant is an induction on
trace length, the local lookup consists of 28 fixed finite cases, the
forward trace construction concatenates blocks of length at most 5, and
the converse uses bounded iteration through $T+1$ parent steps.  Thus the
argument formalizes in PA and proves
\[
 \mathrm{PA}\ \vdash\
 \operatorname{Halt}(R,\mathrm{blank})
 \ \leftrightarrow\
 \operatorname{Halt}(M,\mathrm{blank}).
\]
No assumption of consistency, standard-model soundness, or nonhalting
enters this quotient theorem.


\subsection{The interacting five-step case}
Compatibility is substantive: an apparent seven-edge matching in an
earlier grammar allowed only five folds simultaneously. That grammar was
discarded. The retained grammar passes the simultaneous check, including
a case whose helper itself participates in another fold.

For raw control 2036 reading 1 with its required left-neighbor bit 1,
the raw step writes 1, moves right, and enters 2037. The folded table
executes these five rows; head positions are relative to entry.
\begin{center}\small
\begin{tabular}{@{}rrrrcr@{}}
\toprule
State & Head & Read & Write & Move & Next\\
\midrule
2036 & 0 & 1 & 1 & L & 2046\\
2046 & -1 & 1 & 1 & R & 2043\\
2043 & 0 & 1 & 1 & L & 2005\\
2005 & -1 & 1 & 1 & R & 2054\\
2054 & 0 & 1 & 1 & R & 2037\\
\bottomrule
\end{tabular}
\end{center}
Only the two known cells are read, their contents are unchanged, and
the final state and head exactly match the raw step. This illustrates
why disjoint pairs alone are insufficient: the composed helper behavior
must also be checked.

% Standalone section for the PA195 article. Requires amsmath, amssymb,
% amsthm, hyperref and theorem/lemma/proposition environments in the parent.
% All claims below are mathematical proofs. They are not claims of a newly
% compiled Lean endpoint for the new literal machine.

\section{The PA-provable consistency equivalence}
\label{sec:pa-calibration}

Let $M_*$ denote the delivered 195-state table, with its designated blank
initial configuration, and let $H_*$ be the arithmetic assertion that this
literal machine has a finite halting trace. Throughout this section,
$\operatorname{Con}(\mathrm{PA})$ means that a fixed standard, recursively
checked first-order PA proof system has no proof of $0=1$. Changing between
the usual PA proof codings gives primitive-recursive proof translations
whose correctness PA proves. The conclusion concerns the arithmetic
sentence for the exact table, rather than an external assertion about its
standard execution.

\begin{theorem}\label{thm:pa-biconditional}
For the table $M_*$ in this package,
\[
 \mathrm{PA}\vdash
 \bigl(\neg H_*\leftrightarrow\operatorname{Con}(\mathrm{PA})\bigr).
\]
This is an ordinary mathematical proof of a PA-provability statement.
The new endpoint has not been compiled as an additional Lean or Rocq
theorem. The executable certificates supply the finite literal-program
facts used in the proof; the arguments below explain their arithmetic
meaning and the uniform proof transformations.
\end{theorem}

There are three levels to connect: PA proofs, the numerical proof calculus
searched by the register program, and literal binary-tape traces. The
published PA196 development supplies the PA-refutation compiler for the
original calculus~\cite{chang}. Its other internal direction is not
assumed here. We first justify the incoming diagonal replacement used
by $M_*$~\cite{diagonal}, and then supply the reverse proof
transformation and the finite-trace argument, developing the syntactic
route identified in the earlier bridge notes~\cite{bridge}.

\subsection{The numerical calculus and its diagonal replacement}

Write
\[
 p(a,b)=\frac{(a+b)(a+b+1)}2+b,
 \quad I(A,B)=4p(A,B),\quad M(a,b)=4p(a,b)+1,
\]
\[
 U(a,B)=4p(B,a)+2,\qquad E(a,b)=4p(a,b)+3.
\]
The notational convention is $\neg A=I(A,1)$, conjunction and existential
quantification being their classical abbreviations. Code $0$ is truth,
and code $1=M(0,0)$ is the designated contradiction. The quantifier
constructor has the transposed Cantor arguments displayed above.

The diagonal calculus $\mathcal C_\Gamma$ retains the six PA196 heads,
the dominated propositional head, modus ponens, and closed
generalization. Its guarded gate is
\begin{align*}
 \Gamma(B)&=I(B,U(B,B)),\\
 V(a,b,c)&=I(M(a,b),U(c,M(a,b))),\\
 P(a,b,c)&=I(c,M(c,b)),\\
 S(a,b,c)&=U\bigl(b,
       (M(a,b)\to P(a,b,c))\wedge(M(c,b)\to c)\bigr),\\
 Q(a,b,c)&=U\bigl(c,S(a,b,c)\to\neg P(a,b,c)\bigr),\\
 F_\Gamma(a,b,c)&=\neg\bigl(\Gamma(c)\to(V(a,b,c)\to Q(a,b,c))\bigr).
\end{align*}
The gate is available when $c\ne a,b$; a failed guard emits $0$.
The logical rules and remaining axiom families are exactly those of
\texttt{PA196/Logic/Basic.lean}, with this one replacement of its
\texttt{bundle}.

Every recursive formula child of a nonzero code is smaller than that
code. Every variable label, including a bound label, occurring in a
formula of code $B$ is strictly below $B$. These assertions follow from
the inequalities $a,b\le p(a,b)$ and the four tags, with code zero handled
separately. Consequently $x_B$ is absent from the decoded formula $B$.
The diagonal clause is a valid vacuity instance, although its binder is
chosen from the numerical code of the entire formula.

Classical implicational reasoning extracts from the gate both
$\Gamma(c)$ and the old gate $G(a,b,c)=\neg(V(a,b,c)\to Q(a,b,c))$.
Every diagonal instance $\Gamma(B)$ is available: use $c=B$ and
$a=b=B+1$. Thus all membership-vacuity instances and pointed separation
used in the original proof remain available.

\begin{lemma}[Recovery of the required equality interfaces]
\label{lem:pa-diagonal-interfaces}
The diagonal calculus derives, uniformly in the numerical labels,
\begin{align*}
 E(a,b)&\to\bigl(M(d,a)\leftrightarrow M(d,b)\bigr),\\
 E(a,b)&\to\bigl(M(a,d)\leftrightarrow M(b,d)\bigr),\\
 E(a,b)&\to R(a,b,z)\quad (z\ne a,b),
\end{align*}
where
\[
 R(a,b,z)=U\bigl(z,M(a,z)\leftrightarrow M(b,z)\bigr)
 \wedge U\bigl(z,M(z,a)\leftrightarrow M(z,b)\bigr),
\]
and it derives $\exists x_a R(a,b,z)$ for $z\ne a,b$.
All these derivations are supplied by primitive-recursive proof templates,
with correctness provable in PA.
\end{lemma}

\begin{proof}
The two forward membership rewrites are retained heads. To obtain the
reverse directions, first assume $a\ne b$. Take
$\chi(t)=M(d,t)$ or $M(t,d)$, initially with $d\ne a$, or take the
diagonal family $\chi(t)=M(t,t)$. Put
\[
 C=E(a,b)\to(\chi(b)\to\chi(a)),\qquad
 \varphi=\neg C,\qquad c=\ulcorner\varphi\urcorner.
\]
Here $c$ is the actual numerical code of the displayed formula, so it is
larger than every label in $C$. In particular, all required guards hold.
Apply the old gate at $(a,a,c)$ and abbreviate its $S,P,Q$ by those
letters. We need only its consequence $\neg Q$, not an unproved general
rule for changing bound-variable names.

The second conjunct in $S$ gives
$S\to U(a,M(c,a)\to\varphi)$. Propositionally,
$\varphi$ implies $E(a,b)$, $\chi(b)$, and $\neg\chi(a)$; hence universal
monotonicity gives
\[
 S\to U(a,M(c,a)\to\neg\chi(a)).
\]
The equality atom has disappeared from this quantified consequent.
The unchanged negative-property specialization lemma now gives
\[
 S\to\bigl(M(c,b)\to\neg\chi(b)\bigr).                 \tag{1}
\]
For precision, this specialization uses only the retained witness
$\exists x_a E(b,a)$, direct rewriting under $E(b,a)$, and vacuity of
$M(c,b)\to\neg\chi(b)$ at $a$. The latter follows from
membership vacuity because $a\ne b,c$ and, in the two nondiagonal
families, $a\ne d$. It is exactly the construction
\texttt{negativeSpecialization} in \texttt{PA196/Logic/Reverse.lean};
none of its premises is primitive equality vacuity.

Under $\varphi$, the point condition $P$ gives $M(c,a)$.
Direct element-slot rewriting under $E(a,b)$ then gives $M(c,b)$.
Together with (1), $S$ contradicts the consequence $\chi(b)$ of
$\varphi$. Classical deduction gives
$\varphi\to(S\to\neg P)$. Generalize this theorem at $c$, distribute
the quantifier through implication, and use
$\Gamma(\varphi)=\varphi\to U(c,\varphi)$. The result is
$\varphi\to Q$. Since the gate supplies $\neg Q$, double-negation
elimination gives $C$, as required.

If $a=b$, the reverse rewrite is propositional. When the fixed label
$d=a$, direct rewriting first takes $M(a,b)$ or $M(b,a)$ to $M(b,b)$;
the already proved diagonal case takes $M(b,b)$ to $M(a,a)$.
This covers every label coincidence without a freshness assumption on
the theorem being derived.

For the neighborhood formula choose $z_0=\ulcorner E(a,b)\urcorner$.
The available diagonal clause $E(a,b)\to U(z_0,E(a,b))$ and the two
bidirectional rewrites yield $E(a,b)\to R(a,b,z_0)$.
To change the binder to an arbitrary $z\ne a,b$, only the case
$z\ne z_0$ remains. Let $P(t)$ be either membership-only matrix in
$R$. Bidirectional rewriting gives
$E(z,z_0)\to(P(z_0)\to P(z))$.
The witness $\exists x_{z_0}E(z,z_0)$, universal distribution, and
vacuity of $P(z)$ at $z_0$ give
\[
 U(z_0,P(z_0))\to P(z).
\]
Its antecedent is vacuous at $z$: its atoms contain only
$a,b,z_0$, and membership vacuity extends through its Boolean
connectives and the outer quantifier by the retained commutation head.
Generalization and distribution therefore give
$U(z_0,P(z_0))\to U(z,P(z))$.
This proves the claimed neighborhood implication for every allowed $z$.
Finally the unchanged witness head and existential monotonicity give
the neighborhood witness, reversing its two arguments by the
propositional symmetry of $R$ if necessary.

Every operation in this proof is a fixed syntactic proof template,
composition, numerical code formation, or a displayed decidable
label-equality case split. The choices $c=\ulcorner\neg C\urcorner$ and
$z_0=\ulcorner E(a,b)\urcorner$ use no unbounded search. PA verifies the
templates by the same proof-checker induction as their original
counterparts. This establishes the asserted uniformity.
\end{proof}

The PA196 refutation compiler uses primitive equality through these
membership-rewriting, neighborhood, and witness interfaces. In its
logical and arithmetic certificate copies, the uses of equality vacuity
are in the reverse-rewriting and neighborhood modules just replaced.
The old gate, all other heads, and the command stack interface remain
available. Substituting Lemma~\ref{lem:pa-diagonal-interfaces}'s templates
therefore gives a PA-verifiable refutation compiler for
$\mathcal C_\Gamma$. In particular,
\begin{equation}\label{eq:pa-forward-calculus}
 \mathrm{PA}\vdash
 \exists p\,\operatorname{Prf}_{\mathrm{PA}}(p,\ulcorner0=1\urcorner)
 \ \longrightarrow\
 \exists d\,\operatorname{Prf}_{\Gamma}(d,1).
\end{equation}
This is reuse of the published proof-producing construction with explicit
replacement templates, rather than an appeal to standard-model
equiconsistency. The 191-item AX word still pushes nine items in the
same positions as the original nine-item command, so its stack
compiler indices need no alteration.

\subsection{A uniform translation of the numerical calculus into PA}

We now prove the converse proof-theoretic direction. Work first in the
fixed conservative definitional extension of PA containing
$\operatorname{pow2}(n)=2^n$ and its usual arithmetical definition.
Only fixed primitive-recursive functions are added. A fixed
primitive-recursive elimination of these definitions converts the
resulting proofs to ordinary PA proofs.

Use the bit predicate
\[
 B(s,z)\;\equiv\;
 \exists u,v\,[s=(2u+1)2^z+v\ \wedge\ v<2^z].
\]
PA proves the usual digit lemmas directly by induction and division
with remainder. In particular $B(s,z)\to z<s$, since
$s\ge 2^z\ge z+1$, and hence PA proves $\neg B(x,x)$.

Let $T(c)$ be the arithmetic formula obtained by decoding the source
code $c$:
\begin{align*}
 T(0)&=(0=0),& T(M(a,b))&=B(x_a,x_b),\\
 T(E(a,b))&=(x_a=x_b),& T(U(a,A))&=\forall x_a\,T(A).
\end{align*}
For a nonzero implication code the translation is the corresponding
implication. There is one numerical alias:
$I(0,0)=0$. Thus the uniform constructor theorem is a proof of
\[
 T(I(A,B))\leftrightarrow(T(A)\to T(B)),
\]
not an unconditional assertion of syntactic identity. At $A=B=0$ it
uses the fixed proof $(0=0)\leftrightarrow((0=0)\to(0=0))$.
Every other case follows from the definition of $T$.

Arithmetic auxiliary bound variables are represented with de Bruijn
indices, or equivalently renamed above all source labels before
substitution. The source label $c$ is therefore absent among the free
variables of $T(c)$ even though the expanded bit formula has its own
bound variables. Quantifier formation closes the indicated free source
label, so nested occurrences of the same label have their usual
shadowing semantics. No assumption that all numerical labels are
distinct is made.

\begin{lemma}[Uniform finite bit comprehension]
\label{lem:pa-bit-comprehension}
Given the syntax code of any arithmetic formula $\theta(z,\vec u)$,
there is a primitive-recursive procedure producing a PA proof of
\[
 \forall n\,\exists s<2^n\,\forall z<n\,
       \bigl(B(s,z)\leftrightarrow\theta(z,\vec u)\bigr).
                                                        \tag{2}
\]
PA proves that the procedure produces a valid proof for every
well-formed formula input, with the indicated fresh variables.
\end{lemma}

\begin{proof}
Take induction on $n$ with the existential formula in (2) as induction
formula. At zero choose $s=0$. Given a witness $s<2^n$, split by
excluded middle on $\theta(n,\vec u)$. If it holds use $s+2^n$;
otherwise use $s$. The fixed bit lemmas say that the first $n$ bits
are unchanged, that the new $n$th bit is respectively one or zero,
and that the result is below $2^{n+1}$. This is a fixed finite
first-order proof pattern with one PA induction instance.

On syntax codes the procedure performs capture-avoiding substitution,
forms the displayed induction instance, and concatenates this fixed
proof pattern with the fixed arithmetic lemma proofs. It never
evaluates $\theta$. Its syntax operations and a fresh-bound-variable
choice are primitive recursive. The PA proof of correctness is
induction on the syntax operations, followed by direct verification
that the generated induction line is an instance of the PA induction
schema and that each subsequent line follows by the stated logical
rule. Thus this is a statement about proof checking, with an arbitrary
formula code as input, not a uniform satisfaction or reflection
principle for PA. Full PA induction is used: $\theta$ is not restricted
to a bounded or quantifier-free formula.
\end{proof}

We spell out the gate supplier because label collisions matter here.
Fix $a,b,c$ with $c\ne a,b$, let $\phi=T(c)$, and regard $x_a,x_b$
as the free values before the inner quantifier is entered.
If $a\ne b$, apply (2) with bound $x_a+x_b+1$ and matrix
\[
 \bigl(B(x_a,z)\vee z=x_b\bigr)\wedge\phi[z/x_b].
                                                        \tag{3}
\]
The witness $s$ satisfies, for every $z$,
\[
 B(x_a,z)\to(\phi[z/x_b]\to B(s,z)),\qquad
 B(s,z)\to\phi[z/x_b],
\]
and $\phi\to B(s,x_b)$. To pass from the bounded comprehension
formula to the displayed unbounded assertions, use
$B(x_a,z)\to z<x_a$ and the fact that a number below $2^n$ has no
set bit at an index at least $n$.

If the \emph{labels} $a=b$ coincide, the first antecedent inside
$T(S)$ becomes $B(z,z)$, not $B(x_a,z)$ with $x_a$ held fixed.
It is therefore false. Apply (2) instead with bound $x_b+1$ and
matrix
\[
 (z=x_b)\wedge\phi[z/x_b].                              \tag{4}
\]
This witness has the required second implication and distinguished
point condition, and the first implication is vacuous. Equality
substitution in (4) supplies $\phi$ at the distinguished point.
The substitutions in (3)--(4) are capture avoiding; in particular
the inner binder corresponding to $b$ is renamed when the witness
expression depends on the original free value of $x_b$.

In both cases the result is a PA proof of
$\exists x_c\,(T(S)\wedge T(P))$. Freshness of $c$ in $\phi$
ensures that choosing the set witness does not alter the property.
Consequently PA proves $\neg T(Q)$. The guarded membership-vacuity
clause has a first-order proof, and $T(\Gamma(c))$ has a first-order
proof by the code-freshness lemma. Classical propositional reasoning
therefore proves $T(F_\Gamma(a,b,c))$.
When the guard fails the command emits $0$, with its fixed truth proof.

The source false code is translated to $T(1)=B(x_0,x_0)$.
Its negation is a fixed PA theorem, uniformly under all binders.
Thus source negation, conjunction, and existential abbreviations have
PA-provable equivalences to the ordinary classical connectives, even
if a surrounding source quantifier binds label zero. This deals with
the false-code convention independently of the truth-code alias.

The other head suppliers are fixed first-order proof templates:
the witness-elimination head, quantifier commutation, the two equality
congruence heads, an equality witness, and quantifier distribution.
Coincident labels are handled by capture-avoiding instantiation, with
a numerical equality case split if needed. The dominated
propositional head is a classical tautology; its arithmetic validity
does not require the dominance restriction. Generalization and modus
ponens use their ordinary PA proof constructors, together with the
implication-constructor equivalence above.

\begin{lemma}[Syntactic refutation translation]
\label{lem:pa-refutation-translation}
There is a primitive-recursive function $J$ such that
\[
 \mathrm{PA}\vdash\forall d\,
 \left(\operatorname{Prf}_{\Gamma}(d,1)
 \to\operatorname{Prf}_{\mathrm{PA}}
             (J(d),\ulcorner0=1\urcorner)\right).
                                                        \tag{5}
\]
\end{lemma}

\begin{proof}
Decode the source derivation by recursion on its finite proof code.
At an axiom node insert the appropriate supplier just described.
At modus ponens concatenate the two translated proofs and the
compositional implication proof; at generalization append the
corresponding quantifier rule. At a code-zero alias emit or use the
fixed truth equivalence. Recursion on a malformed proof is made total
by an arbitrary fixed output; the correctness statement has the
displayed proof-validity hypothesis.

PA proves the syntactic invariant
$\operatorname{Prf}_{\Gamma}(d,A)\to
  \operatorname{Prf}_{\mathrm{PA}}(J_0(d),T(A))$
by induction on this coded recursion. Constructor children decrease
strictly; proof concatenation, label closing and fresh-variable
renaming have primitive-recursive definitions and their usual
syntactic correctness inductions. The gate case is precisely
Lemma~\ref{lem:pa-bit-comprehension} and the two label cases above.
Append the fixed theorem $\neg B(x_0,x_0)$ at a proof of code one,
then eliminate the fixed definitional extension. This gives (5).

This argument proves one quantified assertion about proof codes
inside PA. It does not merely produce a PA proof separately for each
standard formula and then assert a uniform conclusion. The inductive
verification of the compiler, including the schema supplier, is the
step that makes (5) applicable to nonstandard codes in every model
of PA.
\end{proof}

\subsection{The exact proof-store invariant}

The numerical command store is $(t,w,d,h)$. Both $t$ and $h$ carry
formula codes; $w$ is the Cantor-coded stack and $d$ is an arithmetic
parameter. Annotate $t,h$ with PA proofs of $T(t),T(h)$, and annotate
each live stack element with the corresponding proof. The clean store
uses the fixed proof of $T(0)$.

Here is a precise convention that avoids treating Cantor pairing as a
free list constructor. Let $\operatorname{hd},\operatorname{tl}$ be
the two Cantor projections. At zero both are zero. For $w>0$,
$\operatorname{tl}(w)<w$: writing $w=p(a,b)$, if $b>0$ then
$p(a,b)>b$, and if $b=0$ the inequality is immediate. Define
\[
 \delta(w)=\min\{k\le w:\operatorname{tl}^{k}(w)=0\}.
\]
PA proves the set nonempty by induction on $w$.
The canonical stack has exactly these $\delta(w)$ heads, and its
proof annotation is a sequence of that length. Its invariant is
\[
 \forall i<\delta(w)\quad
 \operatorname{Prf}_{\mathrm{PA}}
 \bigl(\pi_i,T(\operatorname{hd}(\operatorname{tl}^{i}(w)))\bigr).
                                                        \tag{6}
\]
Bounded minimization and bounded sequence operations make this an
ordinary primitive-recursive invariant.

Pushing a code $a$ onto $w$ prepends its proof, except when $a=w=0$;
in that one case $p(0,0)=0$ and the canonical annotation remains
empty. Popping nonzero $w$ removes the first proof. Popping zero
returns code zero with its fixed truth proof and leaves an empty
annotation. These statements are exact consequences of
$\operatorname{decode}(p(a,b))=(a,b)$ and
$p(a,b)=0\leftrightarrow a=b=0$, and PA proves each of them.

Now inspect the actual command updates.
AX pushes the saved top, a generalization of the old $h$, and the
seven ordinary head instances, in their specified order, and leaves
the guarded diagonal gate in $t$ with $d=h=0$.
The nine pushed entries have the PA proofs provided above.
For MP with offset $x$, put $y=d+x$.
If $I(t,y)=h$, combine the proofs for $t,h$ to obtain a proof for
$y$, push the truth code, and set $(t,d,h)=(y,y,0)$.
If $I(t,y)=h=0$, then $t=y=0$ and the same operation is justified
by the truth proof. On a failed match the command outputs truth,
retains the stack, and sets $(t,d,h)=(0,y,0)$.
Shift, when it does not halt, sets
\[
 (t,w,d,h)\longmapsto
   \bigl(h,\operatorname{tl}(w),d,\operatorname{hd}(w)\bigr),
\]
and moves the annotations accordingly.

The literal Shift test is $t+h=1$. PA proves that this implies
$t=1$ or $h=1$. In either case the relevant annotation, followed by
the fixed proof $\neg T(1)$, yields a PA refutation. In particular
the invariant must track both visible formula registers, rather
than just $t$.

Every annotation update is primitive recursive. The local statement
just proved holds for \emph{every annotated store}: a successful
command preserves the annotation invariant, and a halting command
produces a PA refutation. PA induction therefore propagates the
invariant through an arbitrary coded finite command sequence,
including one starting at an already live scheduler store.
In particular, if execution from the clean store halts then a PA
refutation exists. Under $\operatorname{Con}(\mathrm{PA})$, every
annotated command has a nonhalting output and its successor store
is annotated again. This is the internal replacement for the
upstream external HF-model soundness argument.

\subsection{Why the finite runtime certificates are sufficient in PA}

The source and physical verifications in this package have the
following mathematical content. A frame consists of its source
path, a gap, and the complete register file. A value $v$ is stored
as $0\,1^{v+1}$. The physical gap is represented by an affine run
length $n+c$, with arbitrary $n\ge0$; register contents and the
complete tail are left symbolic. The initial file has at least ten
blocks, and this lower bound is preserved by every primitive:
increment appends one extra block and decrement retains the length.
The greatest source path length is 37, below the header constant 48.

For each increment, each jump, each zero and positive decrement
branch, and the halt site, the delivered finite trace certificate
starts at the exact frame and ends at the exact successor frame
or at the distinguished halt control. There are 4,815 cases for
the 3,700-site source. These are symbolic branch contracts, rather
than 4,815 examples of numerical register inputs.

There are only four kinds of step in their symbolic interpreter:
an actual table row; a change of notation for an identical tape;
a whole run of one self-loop; and a register-file excursion.
Notation changes preserve every concrete instance. A self-loop
over a run of length $an+c$ has the stated effect by induction on
that length. The register excursions are marking a register,
incrementing it, zero decrement, and positive decrement.
Their proofs in \texttt{M6/Native.lean} are inductions on the finite
file and its unary blocks. They use exactly the specified engine
rows and the append-one tail row. These rows are unchanged by the
new decoder and the seven simultaneous folds.

Consequently the soundness proof of
\texttt{M6/Runtime.lean}, specifically \texttt{sstep\_sound} and
\texttt{exec\_sound}, applies to each case. It gives exact
finite reachability for every natural gap and every valid full
register file. The checker only accepts after it has traversed at
least one actual row, so each completed primitive segment has
strictly positive physical length. A nonhalting contract ends at
the root control, not at the halt control.

The same assertion is provable arithmetically, for the following
specific reason. Finite tapes are coded finite words; runs,
concatenation, frame construction, and the register operations are
primitive recursive. The proofs for a self-loop and each of the
four register excursions use induction on a natural run length or
a coded finite list. Their witness traces can be produced by the
same bounded scans and primitive recursions. Repeating the four
symbolic-step cases and concatenating their traces is again a
primitive recursion. PA proves these inductions directly. The
candidate-specific part is a finite list of table equalities and
completed symbolic traces, all decidable closed arithmetic facts.
Thus PA proves the exact, positive source-to-physical segment
assertion uniformly in coded frames. An external theorem about
standard tapes alone would not suffice; the just-described coded
induction is the internalization used here.

The AX source binding checks every primitive and relative jump in
each of its 191 macro blocks against the inherited total arithmetic
macro, and checks every exterior successor. The free-algebra emitter
check supplies its exact numerical output formulas. Hence the
source-level command and scheduler contracts are inherited with
the gate replaced as stated. The scheduler's other 169 sites are
unchanged. Its finite command-word decoding, the Cantor-pairing
routines, and the fixed AX macro composition are primitive-recursive
computations with PA-provable terminating bounds. More specifically,
the construction \texttt{parsed\_command\_step} in
\texttt{PA196/Source/CommandTransfer.lean} takes the next decoded
command at a fetch boundary, executes it, and returns to the next
fetch boundary in a strictly positive number of primitive source
steps whenever the decoded command has a nonhalting output. It applies
to every store and every scheduler pair $(n,l)$ satisfying that explicit
output premise; the annotated-store invariant under consistency
supplies the premise used below.
The source binding retains this construction with the new AX
composition. The store need not be reset between decoded words;
the proof-store invariant is preserved from one command boundary
to the next.

The complete blank startup is checked for exactly 58,154 steps,
including the entire final nonzero support, the head position,
and the state. It ends at the stated clean source frame.
In native labels that state is 84 and the head coordinate is $-2$;
the injective dense relabeling is checked separately. As a fixed
finite computation, the startup equality and the absence of an
earlier halt are PA-provable.

\begin{lemma}[Reflection of a finite physical halt]
\label{lem:pa-positive-reflection}
PA proves that a finite halting trace of $M_*$ yields a PA
refutation.
\end{lemma}

\begin{proof}
Work in PA and assume $\operatorname{Con}(\mathrm{PA})$.
At a command boundary, decode the next command by the fixed
primitive-recursive function
$\operatorname{parse}(\operatorname{seed}(n,l))$.
The proof-store invariant makes its Shift halt test impossible,
and its nonhalting update supplies the next annotated store.
The positive \texttt{parsed\_command\_step} contract just described
then supplies a finite source execution to the next command
boundary. This is exactly the strengthened syntactic version of
the \texttt{SoundBoundary}/\texttt{boundary\_progress} argument in
\texttt{PA196/Source/Nonhalting.lean}; the external appeal to an
HF model is replaced by the PA proof annotations and consistency.

For any given $T$, primitive recursion through $T+1$ such command
boundaries therefore constructs more than $T$ primitive source
steps from the clean initial source frame. Every command segment
has positive length and ends at the next nonhalt fetch boundary.
The intermediate macro routines terminate by their stated
arithmetic contracts. Thus no source halt can occur inside an
unfinished final macro either. Concatenating the exact positive
physical segments for these primitive steps constructs a physical
run from the same initial frame longer than $T$.

Now suppose there were a blank-tape halting trace of length $T$.
If it stops within startup, the fixed startup computation refutes
it. Otherwise concatenate startup and the physical segments just
constructed. Their total length exceeds $T$. Determinism forces
the alleged halted configuration to be the same configuration
at that index in the constructed run. With the partial operational
semantics there is no outgoing row at halt, a contradiction.
With an absorbing totalization of halt the argument is equivalent:
the run could never again reach a nonhalt root frame.

No inverse decoding of an arbitrary intermediate tape is used.
Nor do we need an unbounded least-time search for the next frame.
The chosen number $T+1$ of completed command segments is a
bounded coded iteration. PA proves sequence concatenation and the
determinism comparison by induction on finite trace length.
We have therefore proved in PA that consistency excludes $H_*$,
or equivalently $H_*\to\neg\operatorname{Con}(\mathrm{PA})$.
\end{proof}

For clarity, exact reachability to a nonhalt endpoint already rules
out a premature halt inside a primitive segment: a deterministic
partial run cannot pass through a state with no outgoing row.
The required runtime premise is therefore the exact positive
endpoint contract verified here, not a disjunction allowing an
unexplained early halt instead of the next source frame.

\subsection{Completing the other direction}

By (\ref{eq:pa-forward-calculus}), a PA refutation produces a
$\mathcal C_\Gamma$ refutation. The unchanged nine-item command
compiler turns that derivation into a finite command word which
halts, using the new gate supplier where required. Here one uses
the existing proof compiler's protected-stack-tail and scheduling
identities, which permit the actual live store; no reset between
candidate words is being assumed. The unchanged scheduler reaches
the compiled word, or halts earlier, by the same bounded candidate
induction as in the original compiler. The total arithmetic macro implementations and
the source-to-physical contracts then give a finite blank-tape
halt of the exact delivered table. These are proof-producing
finite computations and their correctness inductions take place
in PA. Therefore
\[
 \mathrm{PA}\vdash
 \neg\operatorname{Con}(\mathrm{PA})\to H_*.
\]
Combining this with Lemma~\ref{lem:pa-positive-reflection} and
classical contraposition proves Theorem~\ref{thm:pa-biconditional}.

\subsection{Logical strength and the role of ConciseIPRA}

The proof does not replace PA by intuitionistic primitive recursive
arithmetic, the subject of the ConciseIPRA project~\cite{ipra}.
The finite-comprehension supplier uses an induction
instance for an arbitrary formula, which is precisely why a short
iPRA proof checker cannot simply be substituted as the searched
theory. The transferable ideas are smaller derivable interfaces,
freshness read from syntax, and uniform proof-certificate
manipulation. Here the diagonal code-freshness device is already
present in the incoming PA research~\cite{diagonal}; the new reduction is the
compatible physical allocation of that source, not a new claim
that iPRA has PA's consistency strength.


\section{Verification, reproduction, and evidence limits}
\label{sec:verification}
\subsection{Reproduction without a search}
From the archive's top-level directory, run
\begin{verbatim}
python3 -B pa195/build_verify.py
python3 -B verification/verify_independent.py pa195 --concrete
python3 -B verification/check_engine.py
\end{verbatim}
Only Python 3 and its standard library are required. No network access,
optimizer, SAT solver, or proof-assistant installation is needed for
these checks. Python optimization mode \code{-O} must not be enabled,
because the preserved source parser uses assertions. To regenerate
derived certificates, add \code{--export} to the first command.

The builder pins the inherited files, reconstructs both tables, checks
the complete state renaming, validates all incoming-history conditions
and local quotient traces, verifies the source graph and free-algebra
emitter identity, checks both universal symbolic contract lists, and
replays the complete startup. The independent verifier imports none of
the compiler or symbolic-runtime code. It parses the source itself,
resolves control through subtree intervals, checks the literal machine
and conjugacy, and performs two independent startup simulations.

\begin{center}\small
\begin{tabular}{@{}p{0.48\linewidth}p{0.44\linewidth}@{}}
\toprule
Evidence & Result\\
\midrule
Literal table reconstruction & 195 states, 390 rows; exact SHA-256\\
Graph reachability & All 195 working controls\\
Local quotient & 22 incoming rows; 28 exhaustive two-cell cases\\
Positive local durations & Ten length-1, seventeen length-3, one length-5\\
Source identity & 3,700 opcodes and all ordinal successors unchanged\\
Emitter identity & All 191 tokens; full stack and guarded top\\
Symbolic primitive verification & All 4,815 cases, for each of two tables\\
Register oracle prerequisites & All 46 protected rows identical\\
Blank startup & 58,154 steps; state, head, and complete tape exact\\
Independent concrete frames & 19,260 runs; 8,234,830 literal steps\\
\bottomrule
\end{tabular}
\end{center}

The concrete frames use gaps 48, 49, and 61, variable register
contents, both decrement branches, and different nonempty spare tails.
Their longest measured run is 708 transitions. These are transcription
and implementation controls; the universal proof is supplied by the
symbolic contracts and their loop lemmas. Similarly, the symbolic
evaluator's maximum of 497 is an \emph{abstract} execution length, not
a bound on the physical time for arbitrary unary register contents.

\subsection{What the packet establishes}
The seven-state quotient is supported by a complete finite all-tape
proof. The source compilation is supported by exact graph identity,
total symbolic contracts, and the inherited register-block proofs.
The PA consistency calibration is given as an ordinary proof with
uniform proof-code transformations, including the truth-code alias,
the zero-stack alias, variable-label collisions, and both directions
of finite trace transport.

No finite prefix of the actual blank run is used as evidence of
nonhalting. The runtime and startup verifications do not themselves
prove PA's consistency. The theorem only transfers consistency to
nonhalting inside PA; it does not assert that PA proves either of
those equivalent statements. No global minimality claim is made for
195 states.

The remaining difference from a fully formalized result is explicit:
the new complete theorem has not been elaborated in Lean or Rocq.
Selected unmodified upstream proof sources are supplied for review,
with exact pinned URLs and hashes, rather than presented as a new
standalone proof-assistant project.

\subsection{Files and provenance}
\file{pa195/} is the deterministic machine packet.
\file{verification/} contains the independent checker and audit receipts.
\file{comparison/} contains the complete archive inventory, the direct PA
directory inventory, canonical comparison manifests, and the 25 normalized
supported 196-state tables. \file{reference_sources/} supplies the selected
upstream proof modules cited in the arithmetic and engine arguments.
The original licenses and source hashes are retained.

The article is compiled with
\begin{verbatim}
cd report
latexmk -pdf -interaction=nonstopmode -halt-on-error pa195_conpa.tex
\end{verbatim}
All numerical table data printed below are generated from the delivered
machine, so the PDF and the machine-readable artifact specify the same
390 transitions.

\appendix
\clearpage
\section{Complete transition table}
\label{app:table}
An entry \(wDq'\) means write \(w\), move \(D\), and enter \(q'\).
State 0 is initial. The symbol \(\mathrm H\) denotes the separate halt
label 195. The 195 working-state entries below specify both transitions.
Read down each panel, then proceed to the next panel on its right.
\begin{center}
\fontsize{9.5}{11.3}\selectfont
\setlength{\tabcolsep}{3pt}
\begin{tabular}{@{}rll@{\hspace{17pt}}rll@{\hspace{17pt}}rll@{\hspace{17pt}}rll@{}}
\toprule
State & Read 0 & Read 1 & State & Read 0 & Read 1 & State & Read 0 & Read 1 & State & Read 0 & Read 1\\
\midrule
0 & \(1\,R\,1\) & \(1\,R\,61\) & 49 & \(0\,R\,16\) & \(1\,R\,49\) & 98 & \(0\,R\,91\) & \(1\,R\,107\) & 147 & \(0\,R\,148\) & \(1\,R\,162\)\\
1 & \(1\,R\,82\) & \(1\,R\,1\) & 50 & \(0\,R\,63\) & \(1\,R\,64\) & 99 & \(0\,L\,131\) & \(0\,L\,85\) & 148 & \(0\,R\,78\) & \(1\,R\,149\)\\
2 & \(0\,L\,2\) & \(1\,L\,3\) & 51 & \(0\,L\,39\) & \(1\,L\,51\) & 100 & \(0\,R\,108\) & \(1\,R\,94\) & 149 & \(0\,R\,150\) & \(1\,R\,157\)\\
3 & \(0\,L\,4\) & \(1\,L\,3\) & 52 & \(0\,R\,22\) & \(1\,R\,52\) & 101 & \(0\,L\,109\) & \(0\,L\,2\) & 150 & \(0\,R\,126\) & \(1\,R\,151\)\\
4 & \(0\,L\,6\) & \(1\,L\,3\) & 53 & \(0\,L\,11\) & \(0\,R\,49\) & 102 & \(0\,R\,38\) & \(1\,R\,68\) & 151 & \(0\,R\,157\) & \(1\,R\,153\)\\
5 & \(1\,R\,7\) & \(0\,L\,5\) & 54 & \(0\,R\,58\) & \(0\,L\,16\) & 103 & \(0\,R\,59\) & \(1\,R\,110\) & 152 & \(0\,L\,165\) & \(1\,R\,78\)\\
6 & \(0\,L\,9\) & \(1\,L\,3\) & 55 & \(0\,R\,65\) & \(1\,R\,66\) & 104 & \(0\,R\,59\) & \(1\,R\,111\) & 153 & \(0\,R\,154\) & \(1\,R\,155\)\\
7 & \(0\,R\,7\) & \(0\,R\,10\) & 56 & \(0\,R\,59\) & \(1\,R\,67\) & 105 & \(0\,R\,66\) & \(0\,L\,16\) & 154 & \(0\,R\,116\) & \(1\,R\,152\)\\
8 & \(0\,R\,14\) & \(1\,R\,13\) & 57 & \(0\,R\,52\) & \(1\,R\,57\) & 106 & \(0\,R\,14\) & \(1\,R\,110\) & 155 & \(0\,R\,156\) & \(1\,R\,136\)\\
9 & \(0\,L\,15\) & \(1\,L\,3\) & 58 & \(0\,R\,48\) & \(0\,R\,25\) & 107 & \(0\,R\,21\) & \(1\,R\,112\) & 156 & \(0\,R\,78\) & \(1\,R\,78\)\\
10 & \(0\,L\,16\) & \(0\,R\,10\) & 59 & \(0\,L\,28\) & \(1\,L\,68\) & 108 & \(0\,R\,113\) & \(1\,L\,20\) & 157 & \(0\,L\,144\) & \(1\,R\,152\)\\
11 & \(0\,L\,26\) & \(1\,R\,18\) & 60 & \(0\,L\,102\) & \(1\,L\,30\) & 109 & \(0\,R\,30\) & \(1\,R\,65\) & 158 & \(0\,R\,159\) & \(1\,R\,160\)\\
12 & \(0\,R\,19\) & \(1\,R\,20\) & 61 & \(0\,R\,0\) & \(1\,R\,61\) & 110 & \(0\,R\,38\) & \(1\,L\,108\) & 159 & \(0\,R\,160\) & \(1\,R\,137\)\\
13 & \(0\,L\,79\) & \(1\,L\,31\) & 62 & \(1\,R\,125\) & \(1\,L\,51\) & 111 & \(0\,R\,121\) & \(0\,L\,16\) & 160 & \(0\,R\,19\) & \(1\,R\,161\)\\
14 & \(0\,L\,11\) & \(0\,R\,22\) & 63 & \(0\,R\,70\) & \(1\,R\,\mathrm{H}\) & 112 & \(0\,R\,114\) & \(0\,L\,16\) & 161 & \(0\,R\,58\) & \(1\,L\,191\)\\
15 & \(0\,L\,23\) & \(1\,L\,3\) & 64 & \(0\,L\,126\) & \(1\,R\,31\) & 113 & \(0\,R\,24\) & \(0\,R\,47\) & 162 & \(0\,R\,180\) & \(1\,R\,166\)\\
16 & \(1\,L\,2\) & \(0\,L\,2\) & 65 & \(0\,R\,40\) & \(1\,R\,60\) & 114 & \(0\,R\,48\) & \(0\,R\,32\) & 163 & \(0\,R\,12\) & \(1\,R\,164\)\\
17 & \(0\,R\,24\) & \(0\,R\,32\) & 66 & \(0\,L\,79\) & \(1\,R\,71\) & 115 & \(0\,R\,70\) & \(1\,R\,86\) & 164 & \(0\,R\,116\) & \(1\,R\,138\)\\
18 & \(0\,R\,84\) & \(0\,L\,95\) & 67 & \(0\,R\,53\) & \(1\,L\,72\) & 116 & \(0\,R\,19\) & \(1\,R\,118\) & 165 & \(0\,R\,137\) & \(1\,R\,136\)\\
19 & \(0\,R\,28\) & \(1\,R\,29\) & 68 & \(0\,R\,73\) & \(1\,R\,74\) & 117 & \(0\,R\,126\) & \(1\,R\,28\) & 166 & \(0\,L\,190\) & \(1\,L\,177\)\\
20 & \(0\,R\,30\) & \(1\,R\,31\) & 69 & \(0\,R\,69\) & \(1\,R\,75\) & 118 & \(0\,R\,122\) & \(1\,R\,107\) & 167 & \(0\,R\,164\) & \(1\,R\,78\)\\
21 & \(0\,R\,24\) & \(0\,R\,57\) & 70 & \(0\,R\,76\) & \(1\,R\,77\) & 119 & \(0\,R\,94\) & \(1\,R\,86\) & 168 & \(0\,L\,155\) & \(1\,L\,176\)\\
22 & \(0\,R\,33\) & \(1\,R\,22\) & 71 & \(0\,R\,79\) & \(1\,L\,110\) & 120 & \(0\,R\,83\) & \(1\,R\,98\) & 169 & \(0\,R\,152\) & \(1\,R\,58\)\\
23 & \(0\,L\,34\) & \(1\,L\,3\) & 72 & \(0\,L\,98\) & \(1\,R\,54\) & 121 & \(0\,R\,87\) & \(1\,R\,87\) & 170 & \(0\,R\,171\) & \(1\,R\,175\)\\
24 & \(0\,L\,35\) & \(0\,R\,1\) & 73 & \(0\,R\,48\) & \(0\,R\,33\) & 122 & \(0\,R\,17\) & \(1\,L\,109\) & 171 & \(0\,R\,163\) & \(1\,R\,172\)\\
25 & \(0\,R\,36\) & \(1\,R\,25\) & 74 & \(0\,R\,81\) & \(0\,L\,41\) & 123 & \(0\,R\,70\) & \(1\,R\,101\) & 172 & \(0\,R\,161\) & \(1\,R\,174\)\\
26 & \(0\,R\,122\) & \(1\,R\,41\) & 75 & \(0\,R\,50\) & \(1\,R\,132\) & 124 & \(0\,R\,96\) & \(1\,R\,108\) & 173 & \(0\,R\,116\) & \(1\,L\,139\)\\
27 & \(0\,R\,38\) & \(0\,L\,39\) & 76 & \(0\,R\,130\) & \(1\,R\,21\) & 125 & \(1\,L\,39\) & \(1\,L\,99\) & 174 & \(0\,R\,155\) & \(1\,R\,58\)\\
28 & \(0\,R\,40\) & \(1\,R\,41\) & 77 & \(0\,R\,101\) & \(1\,R\,90\) & 126 & \(0\,R\,80\) & \(1\,R\,78\) & 175 & \(0\,R\,154\) & \(1\,R\,172\)\\
29 & \(0\,R\,42\) & \(1\,R\,43\) & 78 & \(0\,R\,92\) & \(1\,R\,93\) & 127 & \(0\,R\,120\) & \(1\,R\,133\) & 176 & \(0\,R\,158\) & \(1\,R\,173\)\\
30 & \(0\,R\,44\) & \(1\,R\,45\) & 79 & \(0\,R\,73\) & \(1\,R\,72\) & 128 & \(0\,R\,124\) & \(1\,R\,123\) & 177 & \(0\,R\,178\) & \(1\,R\,182\)\\
31 & \(0\,R\,21\) & \(1\,R\,46\) & 80 & \(0\,R\,91\) & \(1\,R\,119\) & 129 & \(0\,R\,123\) & \(1\,R\,128\) & 178 & \(0\,R\,179\) & \(1\,R\,181\)\\
32 & \(0\,R\,47\) & \(1\,R\,32\) & 81 & \(0\,R\,95\) & \(1\,R\,96\) & 130 & \(0\,R\,44\) & \(1\,R\,90\) & 179 & \(0\,R\,58\) & \(1\,R\,180\)\\
33 & \(0\,R\,32\) & \(1\,R\,33\) & 82 & \(0\,R\,62\) & \(0\,R\,1\) & 131 & \(0\,L\,5\) & \(0\,L\,131\) & 180 & \(0\,L\,170\) & \(1\,R\,165\)\\
34 & \(0\,L\,69\) & \(1\,L\,3\) & 83 & \(0\,R\,27\) & \(1\,R\,97\) & 132 & \(0\,R\,127\) & \(1\,R\,59\) & 181 & \(0\,R\,156\) & \(1\,R\,116\)\\
35 & \(1\,R\,35\) & \(1\,L\,2\) & 84 & \(0\,R\,89\) & \(1\,R\,88\) & 133 & \(0\,R\,11\) & \(1\,R\,37\) & 182 & \(0\,R\,173\) & \(1\,R\,168\)\\
36 & \(0\,R\,49\) & \(1\,R\,36\) & 85 & \(1\,L\,99\) & \(1\,L\,85\) & 134 & \(0\,R\,135\) & \(1\,R\,185\) & 183 & \(0\,R\,166\) & \(1\,R\,169\)\\
37 & \(0\,R\,28\) & \(1\,R\,20\) & 86 & \(0\,R\,19\) & \(1\,R\,100\) & 135 & \(0\,R\,136\) & \(1\,R\,139\) & 184 & \(0\,R\,116\) & \(1\,R\,173\)\\
38 & \(0\,L\,11\) & \(0\,R\,36\) & 87 & \(0\,R\,58\) & \(1\,R\,58\) & 136 & \(0\,R\,12\) & \(1\,R\,137\) & 185 & \(0\,R\,186\) & \(1\,R\,194\)\\
39 & \(0\,L\,5\) & \(1\,L\,51\) & 88 & \(0\,R\,117\) & \(1\,R\,115\) & 137 & \(0\,R\,86\) & \(1\,R\,138\) & 186 & \(0\,R\,187\) & \(1\,R\,190\)\\
40 & \(0\,R\,24\) & \(0\,R\,25\) & 89 & \(0\,L\,26\) & \(1\,R\,134\) & 138 & \(0\,R\,78\) & \(1\,R\,87\) & 187 & \(0\,R\,116\) & \(1\,R\,188\)\\
41 & \(0\,R\,13\) & \(0\,L\,16\) & 90 & \(0\,R\,95\) & \(0\,L\,131\) & 139 & \(0\,R\,140\) & \(1\,R\,183\) & 188 & \(0\,R\,189\) & \(1\,R\,159\)\\
42 & \(0\,R\,53\) & \(1\,R\,54\) & 91 & \(0\,R\,102\) & \(1\,R\,103\) & 140 & \(0\,R\,141\) & \(1\,R\,166\) & 189 & \(0\,R\,168\) & \(1\,R\,78\)\\
43 & \(0\,R\,55\) & \(1\,R\,56\) & 92 & \(0\,R\,67\) & \(1\,R\,56\) & 141 & \(0\,R\,64\) & \(1\,R\,142\) & 190 & \(0\,R\,191\) & \(1\,R\,192\)\\
44 & \(0\,R\,24\) & \(0\,R\,52\) & 93 & \(0\,R\,55\) & \(1\,R\,104\) & 142 & \(0\,R\,143\) & \(1\,R\,177\) & 191 & \(0\,R\,184\) & \(1\,R\,189\)\\
45 & \(0\,R\,58\) & \(1\,R\,59\) & 94 & \(0\,R\,21\) & \(1\,R\,105\) & 143 & \(0\,R\,12\) & \(1\,R\,144\) & 192 & \(0\,R\,167\) & \(1\,R\,193\)\\
46 & \(0\,R\,60\) & \(0\,L\,16\) & 95 & \(0\,R\,21\) & \(0\,L\,10\) & 144 & \(0\,R\,145\) & \(1\,R\,176\) & 193 & \(0\,R\,168\) & \(1\,R\,122\)\\
47 & \(0\,R\,25\) & \(1\,R\,47\) & 96 & \(0\,R\,59\) & \(1\,R\,41\) & 145 & \(0\,R\,86\) & \(1\,R\,146\) & 194 & \(0\,R\,129\) & \(0\,L\,131\)\\
48 & \(0\,L\,35\) & \(0\,R\,61\) & 97 & \(0\,R\,106\) & \(1\,R\,8\) & 146 & \(0\,R\,147\) & \(1\,R\,170\) &  &  & \\
\bottomrule
\end{tabular}
\end{center}


\section{Archive-by-archive comparison}
\label{app:archives}
Every archive below is at the pinned source commit. The longer
machine-readable inventories preserve exact member names, sizes,
digests, and comparison maps. The entries distinguish mathematical
source simplification from a validated physical state reduction.
{\small
\setlength{\tabcolsep}{5pt}
\begin{longtable}{@{}p{0.34\linewidth}p{0.61\linewidth}@{}}
\toprule
Archive & Contribution and status\\
\midrule
\endfirsthead
\toprule
Archive & Contribution and status\\
\midrule
\endhead
\midrule
\endfoot
\bottomrule
\endlastfoot
\file{U208_ConPA_2026-10-29.zip} & Complete U208, original class A: identify old native roles 172 and 211, with helper 217 (or 214 under alternate native labeling). The local simulation is a finite all-tape argument; full PA calibration is inherited and written, not newly kernel exported.\\[0.45em]
\file{U208_ConPA_Conditional_Context_2026-10-29.zip} & Same conditional-history reduction; primary and alternate presentations give the two original U 208 graph classes. Four local blocks have lengths 7,1,3,1.\\[0.45em]
\file{U208_ConPA_Entry_Context_2026-10-29.zip} & Re-expresses the same reduction by entry guards, with primary and compatibility alternatives. Its saving is not additive with the previous packets.\\[0.45em]
\file{U208_ConPA_History_2026-10-29.zip} & Historical-label presentation of the same two classes and same one-state deletion; includes a failed further reduction.\\[0.45em]
\file{u208_research_2026-10-29.zip} & Both original U 208 classes with symbolic runtime and startup evidence; no further state saving.\\[0.45em]
\file{pa208_research_2026-10-29.zip} & Eighteen additional pairwise nonisomorphic U 208 helper variants, each with a local halting-equivalence argument. Expands accepted U 208 corpus to 20 classes; no accepted 207.\\[0.45em]
\file{Below_U208_Research_2026-10-29.zip} & Source-boundary guards, displaced-head quotient lemma, closed-loop pruning, and broader negative local searches. Pruning 3909 to 3904 source occurrences increases the shared-tree cost. No accepted 207.\\[0.45em]
\file{ConPA_below208_research_2026-10-29.zip} & Two-object obstruction to collapsing source/output cover roles, exhaustive restricted axiom-parameter permutation search, two-row and register-invariant searches. No accepted 207.\\[0.45em]
\file{pa_consistency_research_2026-10-29.zip} & Proves zero-guarded source specialization 3909 to 3904 instructions and preserves the ordinal control graph; no smaller physical table.\\[0.45em]
\file{u208_comparison_research_2026-10-29.zip} & Exact comparison, source/encoding research and a fully specified rejected 207 frame witness. Weak/exact schema equivalence is separate from the still-open runtime/calibration replacement.\\[0.45em]
\file{pa_unary_power_cover.zip} & A unary-matrix power-cover theory and written mutual consistency calibration with full PA. Supplies numerical proof-search interface and logical separations; no new literal TM.\\[0.45em]
\file{pa_tm_research_2026-11-13.zip} & New alternative 196 with head order (1,2,3,5,6,4), unchanged 3812-site CFG, maximum source depth 36 rather than 38. Physical Lean baseline rebuild is supplied. R6-pruned 198 passes runtime but lacks calibrated blank startup.\\[0.45em]
\file{pa_tm_extended_research_2026-11-13.zip} & Four distinct one-row closing-CD variants remain196, with constrained symbolic runtime and unchanged blank startup. Joint source invariants and a generalized register executor strengthen search.\\[0.45em]
\file{conpa_tm_comparison_and_search.zip} & Accepted AD-register-swap196. Two-rotation 195 is rejected by 93 primitive-contract failures. Source geometry changes cannot reuse old contextual folds automatically.\\[0.45em]
\file{pa196_research_2026-11-13.zip} & Exact 157-node representation of the original 3812-site source, down from159; seed stays at source path100000. Claimed prospective 195 allocation loses a required decoder helper. Native HA/full-induction theory is a separate research direction.\\[0.45em]
\file{pa_consistency_research.zip} & New 158-node source and an accepted 196 regrouping; the eighth scanner zero consumes the apparent one-state saving. Its seven-zero 195 version misroots on a concrete source frame.\\[0.45em]
\file{pa_consistency_reduction.zip} & Rebracketed196; H-only Shift 196 with 3811 sites and 57922-step startup; integrated witness gate has 3745 sites but best physical table203. Two optimistic 195 allocations fail, one by an exact short cycle on a source frame.\\[0.45em]
\file{pa_tm_comparison_and_reductions.zip} & Two accepted 196 rotations. Equivalent gate reordering shortens AX 201 to 199 macro items and source 3812 to3773, but best physical compilation198. Finite free-algebra/logic certificates and restricted grammar searches are retained.\\[0.45em]
\file{pa_consistency_research_2026-11-13.zip} & Diagonal vacuity B$\to$U(B,B), exploiting formula-code freshness; AX 201 to191, source 3812 to3700. Initial complete implementation199. Gives detailed written reverse PA proof translation; no new full Lean endpoint.\\[0.45em]
\file{pa_consistency_research_2026-11-13 (2).zip} & New twelve-rotation 196 and diagonal198. Also a command calculus halting at successful MP deriving code1; no new binary table for that command change.\\[0.45em]
\file{conpa_tm_comparison_and_reduction_search.zip} & Two experimental 196 MP commands (result in H, omit zero push), with mathematical command arguments but only bounded physical frame evidence. Two rejected Shift layouts actually halt from blank after 43563 steps.\\[0.45em]
\file{conpa_tm_study_2026-11-13.zip} & Combined 196 implements H-only Shift and MP-in-H, with 4928 uniform contracts; an additional rebracketed Combined 196 has full runtime and startup confirmed in this audit; direct-negation 197 has 3556 sites. A 157-node source receives a larger 207-state realization.\\[0.45em]
\file{pa_consistency_comparison_2026_11_13.zip} & First quantifier head split into modal B plus explosion, with finite Hilbert certificates in both directions. Countermodels obstruct deleting either membership rewrite head or replacing the first quantifier head by reflexivity. Twelve source/table combinations retain separate command/calibration obligations.\\[0.45em]
\file{pa_consistency_comparison_and_new196_2026-11-13.zip} & Three accepted diagonal-calculus 196 layouts: diagonal 196 (191 AX items,3700 sites), separate 196 (190 items,3673 sites), scanner\_best 196 (same 3673 sites, improved address layout). Also a runtime/startup-checked scanner 196 on original source and four auxiliary diagonal closing-CD variants with differential runtime proofs, trit invariants, and exact startup. Short-scanner 195 and U208/PA 196 hybrid startup attempts fail stated obligations.\\[0.45em]
\file{pa_below196_research.zip} & Eight equivalent gate spellings, orientation involutions, exact freshness obstruction R(5)=9, smaller pairing macro that increases DAG cost, H-only guard specialization, and broader finite physical searches. No new accepted smaller table.\\[0.45em]
\file{pa_tm_research_2026-11-13 (2).zip} & New zerojump 196 and an auxiliary CD3 parent (fully replayed in this audit): native27/1 changes 0 L40 to1 R21 under the source26 R6=0 invariant; startup58402. Deleting the resulting prospective helper gives a 195 table with wrong source successor and R6 increment, witnessed also on the actual blank execution.\\[0.45em]
\file{PA_consistency_TM_research_2026-11-13.zip} & New source-only zero-run-five witness (3812 sites,159 internal nodes,depth37), combined-negation/MP 3573-site 157-node witness, and parameter-free residual MP. Fixed/protected physical allocations remain rejected; many bounded quotient searches are negative.\\[0.45em]
\end{longtable}
}


\clearpage
\begin{thebibliography}{9}
\bibitem{chang}
James Chang.
\emph{Nonhalting and PA-independence for a 196-state Turing machine}.
Upstream revision \file{407cdc356f6b586635e71ac92be394255f4bfea4}.
\url{https://github.com/jamesyc/BusyBeaver/tree/407cdc356f6b586635e71ac92be394255f4bfea4}.
The imported paper, source, machine, and proof inventory are in ProveIt,
\file{Computability/BusyBeaver/Research/PA/jamesyc-upstream/pa196/}.

\bibitem{comparison}
ProveIt. \emph{PA196 and the incoming U208 construction}.
Commit \file{906367250a7f96c923147b10df392294bced7be5}.
\url{https://github.com/VladimirReshetnikov/ProveIt/blob/906367250a7f96c923147b10df392294bced7be5/Computability/BusyBeaver/Research/PA/pa196-vs-u208/COMPARISON.md}.

\bibitem{bridge}
ProveIt. \emph{PA196: a route to the missing internal implication}.
Same pinned commit.
\url{https://github.com/VladimirReshetnikov/ProveIt/blob/906367250a7f96c923147b10df392294bced7be5/Computability/BusyBeaver/Research/PA/pa196-vs-u208/PA-BRIDGE.md}.

\bibitem{diagonal}
ProveIt incoming research.
\file{pa_consistency_comparison_and_new196_2026-11-13.zip},
especially the 191-token diagonal calculus and \file{axioms/diagonal196/}.
Archive SHA-256
\nolinkurl{6fcb1cf2878eda152beda0b17cc0242c003908fa8adf07c81a2bbf49870f1866}.
\url{https://github.com/VladimirReshetnikov/ProveIt/blob/906367250a7f96c923147b10df392294bced7be5/docs/incoming/pa_consistency_comparison_and_new196_2026-11-13.zip}.

\bibitem{incoming}
ProveIt. All 27 incoming PA research archives at the pinned commit.
\url{https://github.com/VladimirReshetnikov/ProveIt/tree/906367250a7f96c923147b10df392294bced7be5/docs/incoming}.
The accompanying inventory records every inspected archive and member.

\bibitem{zfc}
ProveIt. ZFC small-machine research and inherited register-engine lineage.
\url{https://github.com/VladimirReshetnikov/ProveIt/tree/906367250a7f96c923147b10df392294bced7be5/Computability/BusyBeaver/Research/ZFC}.

\bibitem{ipra}
ProveIt. \emph{ConciseIPRA}.
\url{https://github.com/VladimirReshetnikov/ProveIt/tree/906367250a7f96c923147b10df392294bced7be5/Logic/PrimitiveRecursiveArithmetic/ConciseIPRA}.
\end{thebibliography}
\end{document}
